\emph{Setting.} By Proposition \claimref{correspondence.reality-conjugation-twistor-form-factor}(2)-(4), $\Sigma$ is a Hermitian form on $\CS$ of signature $(2,2)$ for which $\Som$ and $\Spi$ are null, and by (5) $\sigma = \Sigma|_{\Som\times\Spi}$ is a perfect pairing and the real points of the chart of Theorem \claimref{correspondence.incidence-relation-chart} are exactly $A = I_{7}\hat x$ with $\hat x$ $\sigma$-Hermitian, $I_{7}$ the complex unit of the Fock model (which acts as $I_{6}$ up to sign). The chart map is complex linear (Theorem \claimref{correspondence.incidence-relation-chart}(2)), so the complex point $x + I_{7}y$ has $A = I_{7}\hat x + I_{7}(I_{7}\hat y) = I_{7}\hat x - \hat y$.

\emph{(i)} For $Z = A\pi + \pi$ and $Z' = A\pi' + \pi'$, nullity of the halves and hermiticity give $\Sigma(Z, Z') = \sigma(A\pi, \pi') + \overline{\sigma(A\pi', \pi)}$. With $\sigma$ antilinear in its first slot and $\hat x$, $\hat y$ $\sigma$-Hermitian, the terms in $I_{7}\hat x$ give $-I_{7}\sigma(\hat x\pi,\pi') + \overline{-I_{7}\sigma(\hat x\pi',\pi)} = -I_{7}\sigma(\hat x\pi,\pi') + I_{7}\sigma(\hat x\pi,\pi') = 0$, and the terms in $-\hat y$ give $-\sigma(\hat y\pi,\pi') - \overline{\sigma(\hat y\pi',\pi)} = -2\sigma(\hat y\pi,\pi')$. So $\Sigma|_{L} = -2\,\sigma(\hat y\,\cdot,\cdot)$, which depends on $y$ only. In bases of $\Som$, $\Spi$ dual for $\sigma$ the form $\sigma(\hat y\,\cdot,\cdot)$ is a Hermitian matrix $Y$, and by clause (5) of the Proposition its determinant is a positive multiple of $Q(y)$ (the chart norm on real points is $2\det\hat x$); a change of basis of $L$ multiplies the determinant by a positive number, so the sign of $\det(\Sigma|_{L})$ is the sign of $Q(y)$.

\emph{The check computes all of this on the stored construction.} It rebuilds the Fock model, the Chevalley pairing, the real orthonormal basis $a_{1}, a_{2}, b_{1}, \dots, b_{4}$ with $w_{3} = b_{3} + I_{7}b_{4}$, $\kappa = \Gamma C(b_{4})J$ and $\Sigma = \mu B(\kappa\,\cdot,\cdot)$ exactly as the stored Proposition defines them (with $c = \sqrt{2}C$, a positive rescaling, and the stored phase, here $\lambda = 1$, $\mu = 2$), and the stored chart (the graph line in the stored bases $\pi_{1}, \pi_{2}, e_{1}, e_{2}$, with $m(A)$ read through the spinor identification). It checks that $\Sigma$ is Hermitian of rank four with both halves null and $\kappa^{2} < 0$, that the chart reproduces the four basis vectors $e_{0} = a_{2}$, $e_{1} = b_{2}$, $e_{2} = b_{3}$, $e_{3} = b_{4}$ of the Minkowski complement, and that the lines of real points are $\Sigma$-null. For six imaginary parts $y$ (timelike of both orientations, spacelike, null) and two real parts it finds the Gram matrix of $\Sigma$ on the line, in the basis $(A\pi_{1} + \pi_{1}, A\pi_{2} + \pi_{2})$, independent of $x$, with determinant $4Q(y)$ and trace $4y^{0}$.

\emph{(ii)} A Hermitian $2 \times 2$ matrix is definite iff its determinant is positive, then with the sign of its trace; split iff its determinant is negative; and of rank one iff its determinant vanishes and it is nonzero. With determinant $4Q(y)$ and trace $4y^{0}$ this gives the four cases, the definite lines being positive exactly when $y^{0} > 0$. This sign is a property of the stored $\kappa$ and of the orientation of $e_{0}$: replacing $\kappa$ by $-\kappa$ replaces $\Sigma$ by $-\Sigma$, and reversing $e_{0}$ reverses $y^{0}$; either exchanges the two orientations. (A simplified model with $\Sigma = \left(\begin{smallmatrix}0 & 1\\ 1 & 0\end{smallmatrix}\right)$ and the Pauli assignment $\hat y(e_{0}) = +1$ gives the opposite sign, which is why the sign is computed on the stored construction.)

\emph{(iii)} The boost of the $(e_{0}, e_{1})$ plane with rapidity $\lambda$ sends $te_{0} + re_{1}$ to $(t\cosh\lambda + r\sinh\lambda)e_{0} + (t\sinh\lambda + r\cosh\lambda)e_{1}$; at $\lambda = I_{7}\theta$ and $(t, r) = (0, \rho)$ this is $\rho\cos\theta\,e_{1} + I_{7}\rho\sin\theta\,e_{0}$, with $Q(y) = \rho^{2}\sin^{2}\theta$ and $y^{0} = \rho\sin\theta$. The check finds the Gram matrix $2\rho\sin\theta$ times the identity at the Pythagorean angles $(\cos\theta, \sin\theta) = (\pm\tfrac35, \pm\tfrac45)$ and at $\theta = 0, \pi$ (with $\rho = \tfrac32$): positive for $0 < \theta < \pi$, negative for $\pi < \theta < 2\pi$, zero at $\theta \in \{0, \pi\}$, where $x(0) = x_{0}$ and $x(\pi) = -x_{0}$ are real. The orbit is $2\pi$-periodic since $x(\theta)$ depends on $\theta$ through $\cos\theta$ and $\sin\theta$.

\emph{Reading.} The positive Hermitian structure of Theorem \claimref{synthesis.incidence-vacua-thermofield-real-capacity} is not $\mathrm{Spin}(2,4)$-invariant (its clause (iv)); $\Sigma$ is (Proposition \claimref{correspondence.reality-conjugation-twistor-form-factor}(2)). Under $\Sigma$, by (ii), no line of a real point carries a positive state; the incidence line of a real point is a maximal isotropic plane for a form of Witt index two, a real-capacity count for $\Sigma$ in the sense of Proposition \claimref{one-relation.two-capacities-three-consequences}. By (iii) positivity along an imaginary boost holds on exactly one half strip, $0 < \theta < \pi$ for the stored $\kappa$ and imaginary part along $+e_{0}$.
